% GENERATED FILE. Edit canonical lessons, then run npm run generate:books.
% canonical-source-hash: fnv1a64:5991e3ea48874bb7
% canonical-lessons: IT-C421-molo, IT-C421-pontile, IT-C421-imbarcarsi, IT-C421-sbarcare, IT-C421-salpare

\chapter{Al porto}
\label{ch:it-al-porto}
\hlchaptermodality{421}

\begin{chapteropening}
\textbf{By the end of this chapter:} \emph{I can get around a harbour: the pier, the jetty, boarding, landing and setting sail.}

Complete the last lesson of chapter 421: I can get around a harbour: the pier, the jetty, boarding, landing and setting sail.
\end{chapteropening}

\section[il molo]{il molo — the pier, the jetty (of stone)}

\label{lesson:IT-C421-molo}



\begin{quote}
Before the new one: say the Italian for the refund, then the Italian for to refund.
\end{quote}

\subsection*{You'll want to know: il molo}
\textbf{il molo} — \textquotedblleft{}the pier, the jetty (of stone)\textquotedblright{}.

\begin{quote}
Le barche sono legate al molo.
\end{quote}

\begin{quote}
\emph{The boats are tied to the pier.}
\end{quote}

\emph{Il molo} is built of stone and shelters the harbour.

\begin{cousinweb}
\emph{molo} is inherited from Latin \emph{mōlēs} 'mass, heap, rock' (EtymDB), which could also mean a massive stone structure such as a dam or a pier (Lewis \& Short).
\end{cousinweb}

\begin{grammarlens}[title={What you've built}]
One more word for telling, explaining and describing at B1.
\end{grammarlens}

\subsection*{Your turn}
\begin{itemize}
\raggedright
  \item \emph{Say it:} \emph{il molo}
  \item \emph{Say it:} \emph{il molo}, once more
  \item \emph{Say it:} a sentence of your own with \emph{il molo}
  \item \emph{Recall:} say the Italian for the refund, then the Italian for to refund, then say \emph{il molo} again
\end{itemize}

\begin{tcolorbox}[breakable,colback=teal!4,colframe=teal!35!black,title={Before you move on}]
What does il molo mean? (\textquotedblleft{}The pier, the jetty (of stone)\textquotedblright{}.) Say it once more.
\end{tcolorbox}
\section[il pontile]{il pontile — the landing stage, the wooden jetty}

\label{lesson:IT-C421-pontile}



\begin{quote}
Before the new one: say the Italian for to refund, then the Italian for the pier.
\end{quote}

\subsection*{You'll want to know: il pontile}
\textbf{il pontile} — \textquotedblleft{}the landing stage, the wooden jetty\textquotedblright{}.

\begin{quote}
Dal pontile i bambini guardano i pesci.
\end{quote}

\begin{quote}
\emph{From the landing stage the children watch the fish.}
\end{quote}

Built on \emph{ponte}: a bridge that reaches out over the water.

\begin{cousinweb}
\emph{pontile} is built on \emph{ponte} 'bridge', which is inherited from Latin \emph{pōns, pontis} 'bridge' (EtymDB).
\end{cousinweb}

\begin{grammarlens}[title={What you've built}]
One more word for telling, explaining and describing at B1.
\end{grammarlens}

\subsection*{Your turn}
\begin{itemize}
\raggedright
  \item \emph{Say it:} \emph{il pontile}
  \item \emph{Say it:} \emph{il pontile}, once more
  \item \emph{Say it:} a sentence of your own with \emph{il pontile}
  \item \emph{Recall:} say the Italian for to refund, then the Italian for the pier, then say \emph{il pontile} again
\end{itemize}

\begin{tcolorbox}[breakable,colback=teal!4,colframe=teal!35!black,title={Before you move on}]
What does il pontile mean? (\textquotedblleft{}The landing stage, the wooden jetty\textquotedblright{}.) Say it once more.
\end{tcolorbox}
\section[imbarcarsi]{imbarcarsi — to board (a ship or a plane)}

\label{lesson:IT-C421-imbarcarsi}



\begin{quote}
Before the new one: say the Italian for the pier, then the Italian for the landing stage.
\end{quote}

\subsection*{You'll want to know: imbarcarsi}
\textbf{imbarcarsi} — \textquotedblleft{}to board (a ship or a plane)\textquotedblright{}.

\begin{quote}
Ci imbarchiamo alle nove per l'isola.
\end{quote}

\begin{quote}
\emph{We board at nine for the island.}
\end{quote}

Without \emph{si}, \emph{imbarcare} is to take on board: \emph{imbarcare i passeggeri}.

\begin{cousinweb}
\emph{imbarcarsi} is built from \emph{in-} and \emph{barca} 'boat', which is inherited from Late Latin \emph{barca} (Lewis \& Short).
\end{cousinweb}

\begin{grammarlens}[title={What you've built}]
One more word for telling, explaining and describing at B1.
\end{grammarlens}

\subsection*{Your turn}
\begin{itemize}
\raggedright
  \item \emph{Say it:} \emph{imbarcarsi}
  \item \emph{Say it:} \emph{imbarcarsi}, once more
  \item \emph{Say it:} a sentence of your own with \emph{imbarcarsi}
  \item \emph{Recall:} say the Italian for the pier, then the Italian for the landing stage, then say \emph{imbarcarsi} again
\end{itemize}

\begin{tcolorbox}[breakable,colback=teal!4,colframe=teal!35!black,title={Before you move on}]
What does imbarcarsi mean? (\textquotedblleft{}To board (a ship or a plane)\textquotedblright{}.) Say it once more.
\end{tcolorbox}
\section[sbarcare]{sbarcare — to disembark, to get off (a ship)}

\label{lesson:IT-C421-sbarcare}



\begin{quote}
Before the new one: say the Italian for the landing stage, then the Italian for to board.
\end{quote}

\subsection*{You'll want to know: sbarcare}
\textbf{sbarcare} — \textquotedblleft{}to disembark, to get off (a ship)\textquotedblright{}.

\begin{quote}
Siamo sbarcati dopo tre ore di mare.
\end{quote}

\begin{quote}
\emph{We landed after three hours at sea.}
\end{quote}

\emph{Sbarcare il lunario} — to make ends meet.

\begin{cousinweb}
\emph{sbarcare} is built from \emph{s-} 'out of' and \emph{barca} 'boat', which is inherited from Late Latin \emph{barca} (Lewis \& Short).
\end{cousinweb}

\begin{grammarlens}[title={What you've built}]
One more word for telling, explaining and describing at B1.
\end{grammarlens}

\subsection*{Your turn}
\begin{itemize}
\raggedright
  \item \emph{Say it:} \emph{sbarcare}
  \item \emph{Say it:} \emph{sbarcare}, once more
  \item \emph{Say it:} a sentence of your own with \emph{sbarcare}
  \item \emph{Recall:} say the Italian for the landing stage, then the Italian for to board, then say \emph{sbarcare} again
\end{itemize}

\begin{tcolorbox}[breakable,colback=teal!4,colframe=teal!35!black,title={Before you move on}]
What does sbarcare mean? (\textquotedblleft{}To disembark, to get off (a ship)\textquotedblright{}.) Say it once more.
\end{tcolorbox}
\section[salpare]{salpare — to set sail, to weigh anchor}

\label{lesson:IT-C421-salpare}



\begin{quote}
Before the new one: say the Italian for to board, then the Italian for to disembark.
\end{quote}

\subsection*{You'll want to know: salpare}
\textbf{salpare} — \textquotedblleft{}to set sail, to weigh anchor\textquotedblright{}.

\begin{quote}
La nave salpa all'alba.
\end{quote}

\begin{quote}
\emph{The ship sets sail at dawn.}
\end{quote}

\emph{Salpare l'ancora} — to weigh anchor.

\begin{cousinweb}
The origin of \emph{salpare} is not settled: some trace it to a Greek sailors' word, others to Catalan.
\end{cousinweb}

\begin{grammarlens}[title={What you've built}]
One more word for telling, explaining and describing at B1.
\end{grammarlens}

\subsection*{Your turn}
\begin{itemize}
\raggedright
  \item \emph{Say it:} \emph{salpare}
  \item \emph{Say it:} \emph{salpare}, once more
  \item \emph{Say it:} a sentence of your own with \emph{salpare}
  \item \emph{Recall:} say the Italian for to board, then the Italian for to disembark, then say \emph{salpare} again
\end{itemize}

\begin{tcolorbox}[breakable,colback=teal!4,colframe=teal!35!black,title={Before you move on}]
What does salpare mean? (\textquotedblleft{}To set sail, to weigh anchor\textquotedblright{}.) Say it once more.
\end{tcolorbox}
